% Chapter 4 -- Batch Files.  IBM 3.10 Reference nests the batch subcommands
% under "Batch File Commands" with "... Subcommand" titles (ibm-dos.md
% Sec 8.1; AN-01).  The entry list is the shipped directive table of
% os/milton/batch.c batch_classify: REM ECHO GOTO IF FOR SHIFT CALL PAUSE
% (+ :label and @).  Behaviour: os/milton/command.c run_batch (ECHO on at
% entry; @; %0..%9, %VAR%, %%; GOTO "Label not found" stops the file; IF and
% FOR hand their command to the command dispatcher, so a batch subcommand
% after IF or FOR is not recognized -- defect, bead needed, WL-0119; SHIFT
% keeps %0; CALL nests, "Batch nesting too deep"; PAUSE MSG-DOS-0016),
% run_batch_invoke, command_repl (A:\AUTOEXEC.BAT).  The sample files are
% pubs/screens/disk/*.BAT in A:\SAMPLES of the practice diskette.
% GATES: test-batch test-batch-exec test-autoexec test-pubs-screens
\Gates{test-batch, test-batch-exec, test-autoexec, test-pubs-screens}
\chapter{Batch Files}

\section{What a Batch File Is}
\index{batch files}
A \emph{batch file} is a text file of InitechDOS commands with the extension
\file{.BAT}.  When you type its name, InitechDOS carries out the commands in
it one after another, as if you had typed them at the keyboard.  Batch files
save typing for jobs you do often.

A batch file can contain any InitechDOS command or program name, and the
batch subcommands described in this chapter: \cmd{CALL}, \cmd{ECHO},
\cmd{FOR}, \cmd{GOTO}, \cmd{IF}, \cmd{PAUSE}, \cmd{REM} and \cmd{SHIFT}.

\section{Creating and Running a Batch File}
Create a batch file with \cmd{COPY CON}, or with any program that saves
plain text.  Each command goes on a line of its own.

To run a batch file, type its name without the extension, followed by any
parameters, and press \key{Enter}.  InitechDOS looks for it as it looks for
any external command (Chapter~2): in the current directory and then in the
directories named in \texttt{PATH}.  A program file with the same name and
the extension \file{.COM} or \file{.EXE} is found first.

To stop a batch file, press \key{Ctrl-C}.  The batch file ends and the
command prompt is displayed.

\section{The AUTOEXEC.BAT File}
\index{AUTOEXEC.BAT@\texttt{AUTOEXEC.BAT}}
When the system starts, the command processor carries out the batch file
\file{AUTOEXEC.BAT} in the root directory of drive~A, if there is one, before
it displays the first prompt.  Use it for the commands you want carried out
every time, such as setting \texttt{PATH}.  The practice diskette's
\file{AUTOEXEC.BAT} holds two lines:
\begin{TypeIn}
@ECHO OFF
SET PATH=A:\APPS
\end{TypeIn}

\section{Replaceable Parameters}
\index{replaceable parameters}
The words you type after the name of a batch file are its
\emph{parameters}.  Inside the file, \texttt{\%1} stands for the first
parameter, \texttt{\%2} for the second, and so on up to \texttt{\%9};
\texttt{\%0} stands for the name of the batch file itself.  A parameter you
do not give is replaced by nothing.  A batch file receives at most nine
parameters; any more that you type are ignored.

The name of an environment variable between percent signs, as
\texttt{\%PATH\%}, stands for the value of that variable.  To put a single
percent sign in a command, write two: \texttt{\%\%}.

\section{Labels and the @ Character}
A line that begins with a colon, as \texttt{:START}, is a \emph{label}: it
marks a place that \cmd{GOTO} can jump to.  Labels are not carried out.

While echoing is on (see ``ECHO Subcommand''), each command line is
displayed, after the prompt, before it is carried out.  A line that begins
with \texttt{@} is never displayed.  Most batch files begin with
\texttt{@ECHO OFF}, so that only the output of the commands appears.

\section{Batch Subcommands}
The entries that follow describe the batch subcommands.  The examples run
the sample batch files in the directory \file{SAMPLES} of the practice
diskette; each example first displays the file with \cmd{TYPE}.

% CALL -- command.c run_batch BL_CALL (resolve_bat_path, run_batch_invoke,
% reload and continue; not found MSG-DOS-0002); BATCH_NEST_MAX ("Batch
% nesting too deep").
\begin{DosEntry}{CALL}{}{batch}
\Purpose Runs another batch file and then continues with the next line of
the first.
\Format \SyntaxLine{\kw{CALL} \opt{\ph{d}\kw{:}}\opt{\ph{path}}\ph{filename} \opt{\ph{parameters}}}
\Remarks
\Param{filename} names the batch file to run, without the extension.

\Param{parameters} become \texttt{\%1} through \texttt{\%9} of the called
file.

When the called file ends, InitechDOS continues with the line after the
\cmd{CALL}.  A called file may itself use \cmd{CALL}.  If batch files are
nested too deeply, InitechDOS displays \texttt{Batch nesting too deep}.  If
the file is not found, it displays \texttt{Bad command or file name}.

\cmd{CALL} is used only in batch files.
\Example
\file{COVER.BAT} calls \file{GREET.BAT}; a key was pressed at the
\cmd{PAUSE} message:
\Shot{call}
\end{DosEntry}

% ECHO -- batch.c batch_classify (ECHO ON / ECHO OFF / text); command.c
% run_batch (echo_on starts 1; bare ECHO prints the state).
\begin{DosEntry}{ECHO}{}{batch}
\Purpose Turns the display of batch command lines on or off, or displays a
message.
\Format \SyntaxLine{\kw{ECHO} \opt{\kw{ON}\texttt{|}\kw{OFF}\texttt{|}\ph{message}}}
\Remarks
\Param{ON} displays each command line of the batch file, after the prompt,
before it is carried out.  Echoing is on when a batch file starts.

\Param{OFF} stops the display of command lines.  Messages and the output of
commands are still displayed.

\Param{message} is displayed whether echoing is on or off.  Parameters and
variables in it are replaced first.

Typed alone in a batch file, \cmd{ECHO} displays \texttt{ECHO is on} or
\texttt{ECHO is off}.  Put \texttt{@} before \cmd{ECHO OFF} so that the
\cmd{ECHO OFF} line itself is not displayed.
\Example
\file{GREET.BAT} uses \cmd{ECHO} to display two messages built from its
parameters:
\Shot{echobat}
\end{DosEntry}

% FOR -- batch.c batch_for_parse / batch_for_next_token (set separated by
% spaces, commas; no wildcard expansion), batch_for_subst; command.c
% run_batch BL_FOR (each command to the dispatcher).  Batch files only.
\begin{DosEntry}{FOR}{}{batch}
\Purpose Carries out a command once for each item in a set.
\Format \SyntaxLine{\kw{FOR \%\%}\ph{c} \kw{IN} \kw{(}\ph{set}\kw{)} \kw{DO} \ph{command}}
\Remarks
\Param{\%\%c} is the variable: two percent signs and a single character.

\Param{set} is a list of items separated by spaces or commas.  Each item is
used exactly as written.

\Param{command} is an InitechDOS command or program.  It is carried out once
for each item, with \texttt{\%\%}\ph{c} replaced by the item.

\cmd{FOR} is used only in batch files.

\draft{bead needed: in DOS 3.3 an item with wildcards stands for every
matching file, and FOR can be typed at the prompt with \%c; the shipped
build uses each item as written, in batch files only}
An item that contains wildcard characters stands for every matching file
name.
\Example
\Shot{for}
\end{DosEntry}

% GOTO -- batch.c batch_goto_target (leading ':' optional, case-
% insensitive); command.c BL_GOTO ("Label not found", file stops).
\begin{DosEntry}{GOTO}{}{batch}
\Purpose Continues the batch file at the line after a label.
\Format \SyntaxLine{\kw{GOTO} \ph{label}}
\Remarks
\Param{label} is the name of a label line in the batch file, written with
or without its colon, in capital or small letters.

InitechDOS goes to the label line and continues with the line after it.  If
there is no such label, it displays \texttt{Label not found} and the batch
file ends.
\Example
\Shot{goto}
\end{DosEntry}

% IF -- batch.c batch_eval_if (NOT; ERRORLEVEL n true if >= n; EXIST via
% OPEN, one file; string==string exact, case-sensitive); command.c BL_IF
% (the command goes to the command dispatcher: GOTO, SHIFT, CALL, PAUSE
% after IF are not recognized -- defect, bead needed).
\begin{DosEntry}{IF}{}{batch}
\Purpose Carries out a command only when a condition is true.
\Format \SyntaxLine{\kw{IF} \opt{\kw{NOT}} \kw{ERRORLEVEL} \ph{number} \ph{command}}
  \SyntaxLine{\kw{IF} \opt{\kw{NOT}} \ph{string1}\kw{==}\ph{string2} \ph{command}}
  \SyntaxLine{\kw{IF} \opt{\kw{NOT}} \kw{EXIST} \opt{\ph{d}\kw{:}}\opt{\ph{path}}\ph{filename}\opt{\kw{.}\ph{ext}} \ph{command}}
\Remarks
\Param{NOT} reverses the condition: the command is carried out when the
condition is false.

\Param{ERRORLEVEL number} is true when the exit code of the last program
the batch file ran is \ph{number} or greater.

\Param{string1==string2} is true when the two strings are the same,
character for character; capital and small letters differ.  Parameters are
replaced first, so \texttt{IF "\%1"=="" ...} tests whether a parameter was
given.

\Param{EXIST filename} is true when the file exists.  Name one file; do not
use wildcards.

\Param{command} is the InitechDOS command or program to carry out when the
condition is true.

\draft{bead needed: the shipped IF passes its command to the command
processor, so IF ... GOTO, SHIFT, CALL and PAUSE are not recognized
(Bad command or file name); DOS 3.3 carries them out}
The command may also be a batch subcommand, as in
\texttt{IF ERRORLEVEL 1 GOTO ERROR}.
\Example
\file{CHECK.BAT} tests whether the file named by its parameter exists:
\Shot{if}
\file{SAVE.BAT} copies files with \cmd{XCOPY} and tests its exit code; the
second run finds nothing to copy:
\Shot{iferr}
\end{DosEntry}

% PAUSE -- command.c BL_PAUSE (MSG-DOS-0016, waits for one key).
\begin{DosEntry}{PAUSE}{}{batch}
\Purpose Suspends the batch file until a key is pressed.
\Format \SyntaxLine{\kw{PAUSE}}
\Remarks
\cmd{PAUSE} displays \texttt{Press any key to continue . . .} and waits.
Press any key to continue the batch file, or \key{Ctrl-C} to end it.

Use \cmd{PAUSE} before a step that needs the operator, such as putting
paper in the printer; display the instruction with \cmd{ECHO} first.
\Example
See the example of \cmd{CALL}: \file{COVER.BAT} displays an instruction
and pauses for a key.
\end{DosEntry}

% REM -- batch.c BL_REM (never carried out, never displayed).
\begin{DosEntry}{REM}{Remark}{batch}
\Purpose Puts a remark in a batch file.
\Format \SyntaxLine{\kw{REM} \opt{\ph{remark}}}
\Remarks
\Param{remark} is any text.  InitechDOS does not carry out or display a
\cmd{REM} line.  Use remarks to say what a batch file is for.
\Example
The second line of \file{GREET.BAT}, in the example of the \cmd{ECHO}
subcommand, is a remark.
\end{DosEntry}

% SHIFT -- command.c BL_SHIFT (%1 dropped, %2..%9 move down, %0 kept).
\begin{DosEntry}{SHIFT}{}{batch}
\Purpose Moves the replaceable parameters one place to the left.
\Format \SyntaxLine{\kw{SHIFT}}
\Remarks
After \cmd{SHIFT}, \texttt{\%1} holds what was the second parameter,
\texttt{\%2} the third, and so on; the first parameter is discarded.  A
batch file can use \cmd{SHIFT} to work through its parameters one at a
time, always looking at \texttt{\%1}.  \texttt{\%0}, the name of the batch file, is not changed.
\Example
\Shot{shift}
\end{DosEntry}
